\subsection{The sequencer and the matching engine together}
\label{sec:ha:disabled:seq_and_me}

Both processes die, and both are restarted. This is not the two preceding
subsections happening at once: each of them relies on a component that is, in
this case, also gone.

\subsubsection{What state it holds}
\label{sec:ha:disabled:seq_and_me:state}

Between them, everything the venue knows about work in progress. The sequencer
holds the account of what has been taken; the matching engine holds which orders
are open. Neither holds the other's, and neither can reconstruct it.

\subsubsection{Where the copy that outlives it lives}
\label{sec:ha:disabled:seq_and_me:durable}

The log, and nothing else. What is on disk is the record of what the venue took
and in which order. Nothing on disk says which orders were open.

The difference from the single failures is not that more is lost, but that the
survivor which made each of them recoverable is absent. A matching engine
recovers by asking the sequencer how far it is behind; a sequencer is reminded
of its sessions by components that are still running. When both have gone, the
first has nobody to ask, and what the second must be told is held by the
component that is asking.

\subsubsection{How a restarted instance becomes current}
\label{sec:ha:disabled:seq_and_me:recovery}

An order between the two --- committed, and never given to the matching engine
--- is the case this subsection is about. It is in the log, so the venue took
it. The engine does not hold it, so nothing will answer it. And with both components restarted,
nothing records that it is still to be answered: the sequencer's record that it
was waiting did not survive, and the engine holds nothing to compare against.

Recovery therefore acquires an order it does not otherwise have. The engine
cannot become current before the sequencer is current, because it becomes
current by asking. Where only one of the two failed, that ordering was supplied
by the survivor being already up.

\begin{gap}{BUG-0064}
  An order committed before both components died is applied by neither, answered
  by neither, and recorded as outstanding by neither. It remains in the log
  indefinitely.
\end{gap}

\subsubsection{What the member is told}
\label{sec:ha:disabled:seq_and_me:member}

\begin{requirement}{R-0015}{Trading does not resume while committed orders remain unresolved}
  The venue shall not accept new orders after such a failure until every order
  it had committed has either been given to the matching engine or had an
  execution report produced for it. Whether a member is present to receive that
  report does not bear on this: a member that is slow to return, or that never
  returns, cannot hold the venue shut.
  \rationale{Accepting new orders while earlier ones remain unresolved conceals
  the loss rather than recovering from it: the venue looks healthy to a member
  whose earlier order will never be mentioned again.}
  \uncovered{no scenario restarts both components together}
\end{requirement}

Where the automatic paths cannot resolve a failure, the recourse is a person
rather than a longer timer. That requires the venue to have a state it can be
put into, which is a capability rather than a reaction: a venue needs to stop
trading for reasons that have nothing to do with any failure --- a scheduled
pause, a market-wide event, an instrument in disorderly conditions.

\begin{requirement}{R-0023}{Trading can be halted by a person, and members are told}
  An operator shall be able to halt trading. Members shall be told that the
  venue has halted and shall be told when it resumes, and the halt shall not
  lift of its own accord.
  \rationale{A member that is connected but idle learns nothing from the
  rejection of an order it has not sent, and a member reconnecting into a halt
  would otherwise discover it by trial. Requiring a person to lift it is what
  distinguishes a halt from the automatic, self-clearing states beneath it.}
  \uncovered{the venue has no halted state}
\end{requirement}

\begin{requirement}{R-0024}{A halt that cancels open orders reports every cancellation}
  Where a halt cancels the venue's open orders, an execution report shall be
  sent to the member that placed each one. No order shall cease to be open
  without the member that placed it being told.
  \rationale{Cancelling the open orders is a normal thing for a halt to do, and
  has been done at real venues when a market became disorderly. What makes it
  safe is that each member's view converges on the venue's, which requires the
  venue to say what it cancelled rather than to stop mentioning it.}
  \uncovered{the venue has no halted state}
\end{requirement}

\begin{requirement}{R-0017}{Recovery does not depend on the order the components return in}
  The venue shall reach the same state whichever of the sequencer and the
  matching engine starts first after both have failed.
  \rationale{A launcher restarts what it starts and coordinates with nothing, so
  the order in which they return is not something the venue decides. A component that cannot yet be told
  what it needs must wait to be told, rather than proceed on the assumption that
  there is nothing to hear.}
  \uncovered{no scenario restarts both components together}
\end{requirement}
